\section{The full arithmetic target at an odd coefficient prime}
\label{sec:odd-galois-target}

Fix an odd prime \(p\).  All cohomology in this section is
continuous cohomology with the trivial discrete coefficient
module \(\mathbf F_p\), unless a twist is displayed.  Put
\[
 K_2=\mathbf Q(\mu_\infty,2^{1/\infty}),\quad
 A=\operatorname{Gal}(K_2/\mathbf Q),\quad
 U=\widehat{\mathbf Z}^{\times},\quad
 Y_p=H^1(G_{K_2},\mathbf F_p)^{G_{\mathbf Q}}.
\]
The established affine description gives
\begin{equation}\label{eq:odd-target-affine-extension}
 1\longrightarrow 2\widehat{\mathbf Z}(1)
       \longrightarrow A\longrightarrow U\longrightarrow1,
\end{equation}
with conjugation on the translation kernel given by
multiplication by the corresponding unit.

Let \(\mathcal S_p\) be the set of rational primes
\(\ell\equiv1\pmod p\).  Put
\(v_\ell=v_p(\ell-1)\), and choose an isomorphism from
the maximal \(p\)-quotient of \(\mathbf Z_\ell^\times\)
to \(\mathbf Z/p^{v_\ell}\mathbf Z\).
Its reduction modulo \(p\) defines \(\gamma_\ell\).
Choose also a generator
\(\gamma_0\in\operatorname{Hom}(1+p\mathbf Z_p,\mathbf F_p)\),
extended to \(\mathbf Z_p^\times\) by zero on its
prime-to-\(p\) torsion.  Regard all these as characters
of \(U\) by the respective coordinate projections.

Define \(\delta_\ell\in H^2(U,\mathbf F_p)\) by pulling
back the central extension
\begin{equation}\label{eq:odd-target-cyclic-extension}
 0\longrightarrow\mathbf F_p
   \xrightarrow{\ a\mapsto p^{v_\ell}a\ }
   \mathbf Z/p^{v_\ell+1}\mathbf Z
   \longrightarrow\mathbf Z/p^{v_\ell}\mathbf Z
   \longrightarrow0.
\end{equation}
This normalization uses the chosen cyclic coordinate.
When \(v_\ell>1\), it must not be confused with applying
the ordinary length-two Bockstein directly to
\(\gamma_\ell:G\to\mathbf F_p\).

\begin{theorem}\label{thm:odd-full-galois-target}
Inflation gives an isomorphism
\[
                         H^2(U,\mathbf F_p)\simeq H^2(A,\mathbf F_p).
\]
In fact the corresponding inflation is an isomorphism in
every degree.  The arithmetic inflation images of the
\(\delta_\ell\), for \(\ell\in\mathcal S_p\), form a
basis of \(H^2(G_{\mathbf Q},\mathbf F_p)\).

Order \(\{0\}\cup\mathcal S_p\) with zero first.  For
each \(i<j\), there is a unique finite linear combination
\(D_{ij}\) of the \(\delta_\ell\) having the same
arithmetic image as \(\gamma_i\smile\gamma_j\).
The transgression is an isomorphism
\[
 Y_p\simeq
 \ker\bigl(H^2(A,\mathbf F_p)
                  \longrightarrow H^2(G_{\mathbf Q},\mathbf F_p)\bigr),
\]
and the inverse images of
\(\gamma_i\smile\gamma_j-D_{ij}\), identified with
classes of \(A\) by inflation, form a basis of \(Y_p\).
Every linear combination here has finite support.
\end{theorem}

There is also an explicit two-place formula for the
correction.  Set
\[
              a_\ell=\frac{\ell-1}{p^{v_\ell}}\bmod p
                                            \in\mathbf F_p^\times.
\]
For distinct positive indices \(\ell<m\) one has
\begin{equation}\label{eq:odd-target-two-place-correction}
 D_{\ell m}=
       a_\ell^{-1}\gamma_m(\ell)\,\delta_\ell
          -a_m^{-1}\gamma_\ell(m)\,\delta_m,
 \qquad
 D_{0\ell}=-a_\ell^{-1}\gamma_0(\ell)\,\delta_\ell.
\end{equation}
The evaluations use \(\ell\in\mathbf Z_m^\times\),
\(m\in\mathbf Z_\ell^\times\), and
\(\ell\in\mathbf Z_p^\times\), respectively.
Frobenius is normalized to act on residue fields by the
positive power map.  Formula
\eqref{eq:odd-target-two-place-correction} corresponds
to the extension convention
\eqref{eq:odd-target-cyclic-extension}.

\subsection{Cohomology of the unit and affine groups}

The maximal pro-\(p\) quotient of \(U\) is
\[
            U^{(p)}\simeq\mathbf Z_p
                    \times\prod_{\ell\in\mathcal S_p}
                                  C_{p^{v_\ell}}.
\]
The kernel has no finite quotient of order divisible by
\(p\), and hence has no positive-degree
\(\mathbf F_p\)-cohomology: for its finite quotients,
averaging makes invariants exact, and one passes to the
continuous direct limit.  Thus \(U\) and \(U^{(p)}\)
have the same cohomology with these coefficients.

The cyclic resolution gives one generator in degrees
one and two for each finite cyclic \(p\)-group.  The
degree-two generator is its nonsplit cyclic-extension
class.  The square of a degree-one class is zero, since
\(p\) is odd.  For \(\mathbf Z_p\), continuous
cohomology has only degrees zero and one: on passing
from \(C_{p^n}\) to \(C_{p^{n+1}}\), the degree-two
extension class inflates to zero because the quotient
map lifts to \(C_{p^{n+1}}\) itself; cyclic periodicity
then kills every degree at least two in the direct
limit.  The cyclic calculation and the continuous
direct-limit assertion are recalled in
\cite[II, Proposition~3.4, p.~81; Proposition~4.2,
pp.~87--88]{MilneCFT}.

Tensoring the cyclic resolutions for a finite product,
or the usual Kunneth formula over \(\mathbf F_p\), now
gives
\begin{align}
 H^1(U,\mathbf F_p)
   &=\bigoplus_{i\in\{0\}\cup\mathcal S_p}\mathbf F_p\gamma_i,
       \label{eq:odd-target-unit-H1}\\
 H^2(U,\mathbf F_p)
   &=\left(\bigoplus_{\ell\in\mathcal S_p}\mathbf F_p\delta_\ell\right)
       \oplus
       \left(\bigoplus_{i<j}\mathbf F_p
                                  (\gamma_i\smile\gamma_j)\right).
       \label{eq:odd-target-unit-H2}
\end{align}
These are direct sums.  A continuous cochain on a compact
profinite group, with finite discrete values, factors
through a finite quotient.  An open subgroup of the
displayed product contains every factor outside a
finite set.  Thus a cohomology class involves only
finitely many cyclic coordinates and a finite quotient
of the \(\mathbf Z_p\) coordinate.  This proves the
finite-support assertion in these formulas.

Let \(T=2\widehat{\mathbf Z}(1)\) in
\eqref{eq:odd-target-affine-extension}.  Its cohomology
with trivial coefficients is
\[
 H^0(T,\mathbf F_p)=\mathbf F_p,\quad
 H^1(T,\mathbf F_p)=\mathbf F_p(-1),\quad
 H^b(T,\mathbf F_p)=0\quad(b\ge2).
\]
The action in the middle expression is the inverse
cyclotomic action, since a unit acts on translations
by multiplication.  The central order-two subgroup
\(\{1,-1\}\subset U\) acts by \(-1\) on this
one-dimensional module.  Its invariants are zero,
and averaging kills its positive cohomology.  The
Hochschild--Serre spectral sequence for that subgroup
therefore gives
\[
                         H^a(U,\mathbf F_p(-1))=0
                                      \quad\hbox{for every }a\ge0.
\]
In the spectral sequence of
\eqref{eq:odd-target-affine-extension}, only the
row \(b=0\) remains.  Its edge map is inflation,
which proves the asserted isomorphism in all degrees.
The spectral-sequence and five-term conventions used
here are those of
\cite[II, Remark~1.35, p.~72; \S4, pp.~87--88]{MilneCFT}.

\subsection{The global second cohomology, including its support}

We derive the required global statement directly from
the Brauer local-global theorem.  Put
\(L=\mathbf Q(\mu_p)\),
\(\Delta=\operatorname{Gal}(L/\mathbf Q)\), and let
\(\chi:\Delta\to\mathbf F_p^\times\) be the
cyclotomic character.  Since \(|\Delta|=p-1\),
restriction gives
\[
 H^2(G_{\mathbf Q},\mathbf F_p)
                     \simeq H^2(G_L,\mathbf F_p)^\Delta.
\]
Choose \(\zeta_p\) to identify the trivial module over
\(L\) with \(\mu_p\).  Kummer theory, Hilbert 90,
and the cohomological interpretation of the Brauer
group identify \(H^2(G_L,\mu_p)\) with
\(\operatorname{Br}(L)[p]\); see
\cite[II, Proposition~1.22; IV, Theorem~3.14,
p.~137]{MilneCFT}.
The identification by \(\zeta_p\) is not
\(\Delta\)-equivariant: it twists the action by
\(\chi\).  Consequently
\begin{equation}\label{eq:odd-target-brauer-eigenspace}
 H^2(G_{\mathbf Q},\mathbf F_p)
              \simeq\operatorname{Br}(L)[p]^{\chi},
\end{equation}
where \(\delta b=\chi(\delta)b\) defines the
eigenspace on the right.  Indeed if \(c\) has trivial
\(\Delta\)-action before changing coefficients,
then applying \(\delta\) to \(\zeta_p^c\) multiplies
its additive exponent by \(\chi(\delta)\).

The Brauer fundamental exact sequence
\cite[VIII, Theorem~4.2, p.~241]{MilneCFT}, together
with the local invariant isomorphisms, gives
\begin{equation}\label{eq:odd-target-brauer-exact}
 0\longrightarrow\operatorname{Br}(L)[p]
       \longrightarrow\bigoplus_{w\ {
m finite}}\mathbf F_p
       \xrightarrow{\ \sum\ }\mathbf F_p\longrightarrow0.
\end{equation}
The field \(L\) is totally imaginary.  The group
\(\Delta\) permutes the local summands and acts
trivially on the final sum.  Taking a prescribed
\(\Delta\)-eigenspace is exact, because \(p\) does
not divide \(|\Delta|\).  The nontrivial character
\(\chi\) has zero eigenspace in the final term.
It follows that the \(\chi\)-eigenspace of
\(\operatorname{Br}(L)[p]\) is the direct sum of
the \(\chi\)-eigenspaces of the local permutation
modules.

For a rational prime \(r\) and \(w\mid r\), that
permutation module is \(\mathbf F_p[\Delta/\Delta_w]\).
Its \(\chi\)-eigenspace has dimension one if
\(\chi|_{\Delta_w}=1\), and zero otherwise.
The character \(\chi\) is faithful, so the former
condition is exactly that \(r\) split completely
in \(L\), namely \(r\in\mathcal S_p\).

The same coefficient-change argument over each local
field shows that \(H^2(G_{\mathbf Q_r},\mathbf F_p)\)
is one-dimensional for \(r\in\mathcal S_p\), and
zero otherwise.  At the prime \(p\), the local
decomposition group is all of \(\Delta\), so it
contributes zero.  Real places contribute zero since
their group has order two.  The local invariant
of a Brauer class is unchanged by a field
automorphism, which is the only local equivariance
used here.  Hence localization itself is an
isomorphism
\begin{equation}\label{eq:odd-target-localization-isomorphism}
 H^2(G_{\mathbf Q},\mathbf F_p)
          \xrightarrow{\ \sim\ }
       \bigoplus_{\ell\in\mathcal S_p}
                            H^2(G_{\mathbf Q_\ell},\mathbf F_p).
\end{equation}
In particular the support of every global class is
finite.  There is no residual global sum relation:
the sum target in \eqref{eq:odd-target-brauer-exact}
has trivial character, whereas the coefficient
module here selects the nontrivial \(\chi\)-component.

\subsection{The cyclic classes give the local basis vectors}

For \(r\ne\ell\), the local restriction of the
cyclotomic \(\ell\)-coordinate is unramified.
The restriction of \(\delta_\ell\) therefore factors
through \(G_{\mathbf Q_r}/I_r\simeq\widehat{\mathbf Z}\),
whose second \(\mathbf F_p\)-cohomology is zero.
Thus its only possible support is at \(\ell\).

At \(\ell\), its restriction is nonzero.  To see
this directly, write \(v=v_\ell\).  Local class
field theory identifies the maximal abelian
pro-\(p\) quotient of \(G_{\mathbf Q_\ell}\) with
\[
                             \mathbf Z_p\times C_{p^v}.
\]
The two factors come from the uniformizer and the
\(p\)-primary torsion in \(\mathbf Z_\ell^\times\).
The cyclotomic \(\ell\)-coordinate restricts to an
isomorphism on that torsion factor.  This is the
usual unit action in local cyclotomic reciprocity
\cite[I, Theorem~1.1, p.~20]{MilneCFT}; it also follows
directly from \cite[p.~380, formula~(2)]{LubinTate1965}.
If the restricted extension class were zero, this
character into \(C_{p^v}\) would lift to a
homomorphism into \(C_{p^{v+1}}\).
But every homomorphism from the torsion factor
\(C_{p^v}\) into \(C_{p^{v+1}}\) has image in
\(pC_{p^{v+1}}\), whose image modulo \(p^v\)
cannot generate \(C_{p^v}\).  This is impossible.
The restriction is therefore nonzero, and is a
basis of its one-dimensional local cohomology.

Together with
\eqref{eq:odd-target-localization-isomorphism},
this proves that the arithmetic images of the
\(\delta_\ell\) form the claimed global basis.

\subsection{The exact two-place correction}

Let \(\eta_\ell:G_{\mathbf Q_\ell}\to\mathbf F_p\)
be the unramified character taking arithmetic
Frobenius to one.  With the normalization
\eqref{eq:odd-target-cyclic-extension}, one has
\begin{equation}\label{eq:odd-target-local-cyclic-cup}
 \operatorname{loc}_\ell(\delta_\ell)
                   =a_\ell\,\gamma_\ell\smile\eta_\ell.
\end{equation}
We check the sign and the exponent.  Since
\(\mu_{p^v}\subset\mathbf Q_\ell\), write
\[
       \chi_p(g)=1+p^v\tau(g)\pmod{p^{v+1}},
       \qquad \tau=a_\ell\eta_\ell.
\]
Choose the lift of \(\gamma_\ell\) to the cyclic
character with values in \(\mathbf Z/p^v\mathbf Z\),
and choose a cochain lift of it to
\(\mathbf Z/p^{v+1}\mathbf Z\).
The connecting cocycle for trivial coefficients is
\(\delta_\ell\).  If instead the middle coefficient
module has cyclotomic action, its differential
acquires the additional term
\(\tau(g)\gamma_\ell(h)\).
The latter connecting homomorphism is zero: under
the root-of-unity identifications its coefficient
sequence is
\[
                  1\longrightarrow\mu_p
                   \longrightarrow\mu_{p^{v+1}}
                     \xrightarrow{\ z\mapsto z^p\ }\mu_{p^v}
                       \longrightarrow1,
\]
and Kummer theory identifies its map on first
cohomology with the surjection
\(\mathbf Q_\ell^\times/(\mathbf Q_\ell^\times)^{p^{v+1}}
\to\mathbf Q_\ell^\times/(\mathbf Q_\ell^\times)^{p^v}\).
Therefore
\[
       \operatorname{loc}_\ell(\delta_\ell)
             =-\tau\smile\gamma_\ell
             =\gamma_\ell\smile\tau,
\]
where the last equality is graded commutativity
at the odd prime \(p\).  This proves
\eqref{eq:odd-target-local-cyclic-cup}.

For distinct \(\ell,m\in\mathcal S_p\), the
character \(\gamma_m\) at \(\ell\) is unramified,
with Frobenius value \(\gamma_m(\ell)\).  Thus
the local contribution of
\(\gamma_\ell\smile\gamma_m\) at \(\ell\)
is \(\gamma_m(\ell)\gamma_\ell\smile\eta_\ell\).
At \(m\) its contribution is
\(-\gamma_\ell(m)\gamma_m\smile\eta_m\).
All other local contributions vanish: both
characters are unramified there, or the local
second cohomology is zero.  The same argument for
\(\gamma_0\smile\gamma_\ell\) leaves only the
contribution \(-\gamma_0(\ell)
\gamma_\ell\smile\eta_\ell\) at \(\ell\).
Using \eqref{eq:odd-target-local-cyclic-cup} and
the localization isomorphism proves
\eqref{eq:odd-target-two-place-correction}.
It proves uniqueness as well, because the cyclic
classes have independent local supports.

\subsection{Transgression and the basis of the target}

Apply the five-term sequence to
\(1\to G_{K_2}\to G_{\mathbf Q}\to A\to1\):
\[
 0\to H^1(A,\mathbf F_p)\to H^1(G_{\mathbf Q},\mathbf F_p)
       \to Y_p\xrightarrow{\ \operatorname{tr}\ }
           H^2(A,\mathbf F_p)\to H^2(G_{\mathbf Q},\mathbf F_p).
\]
Kronecker--Weber identifies
\(G_{\mathbf Q}^{\rm ab}\) with \(U\)
\cite[I, Theorem~4.16, pp.~49--50]{MilneCFT}.
The already proved degree-one inflation isomorphism
therefore makes the first map in this sequence
an isomorphism.  Transgression is consequently
injective and its image is exactly the kernel
of arithmetic inflation in degree two.

By \eqref{eq:odd-target-unit-H2}, every class in
\(H^2(U,\mathbf F_p)\) has a unique finite
decomposition into cyclic classes and pairwise
cups.  The cyclic summand maps isomorphically
onto \(H^2(G_{\mathbf Q},\mathbf F_p)\).
Subtracting \(D_{ij}\) from each cup gives
a basis for the inflation kernel.  Injectivity
of transgression supplies its unique inverse
basis in \(Y_p\), proving the theorem.

This determines the full arithmetic target.
It makes no assertion that a target basis vector
is the restriction of a character of full GT.
